\Subhead{1.7.1}{Параметрическое усиление и параметрическая генерация. Параметрические усилители и генераторы}
\textbf{Устройство и энергия.} Варикап --- нелинейная ёмкость: сильная накачка $U_p(t)$ меняет её дифференциальное значение $C_d(t)=\dd q/\dd U|_{U_p(t)}$. Для слабого сигнала $q_s=C_d(t)u_s$. Произведение $\cos\omega_pt\cos\omega_st$ создаёт частоты $\omega_p\pm\omega_s$. Резонансные цепи выбирают нужные составляющие и отделяют мощную накачку от сигнала.

\begin{center}
\begin{circuitikz}[scale=.85,transform shape]
\draw (0,0) to[I,l=$I_s$] (0,2) -- (3,2) to[vC,l=$C_c(U_p)$] (5,2) -- (8,2);
\draw (0,0)--(8,0);
\draw (1,2) to[L,l=$L_1$] (1,0);
\draw (2,2) to[C,l=$C_1$] (2,0);
\draw (3,2) to[R,l=$R_1$] (3,0);
\draw (5,2) to[C,l=$C_2$] (5,0);
\draw (6.5,2) to[L,l=$L_2$] (6.5,0);
\draw (8,2) to[R,l=$R_2$] (8,0);
\draw[dashed,->] (4,3.1)--(4,2.4);
\node at (4,3.45) {Накачка $\omega_p\simeq\omega_s+\omega_i$};
\node at (1.5,-.55) {Сигнальный контур};
\node at (6.4,-.55) {Холостой контур};
\end{circuitikz}
\end{center}
Это двухконтурная схема с нелинейной связующей ёмкостью, перерисованная по конспекту Балакшия (PDF-с.~9, лист~7). $R_j$ представляют потери и нагрузку, внешний ток вводит сигнал. Пунктир обозначает управление ёмкостью накачкой, а не отдельную токовую ветвь; монтажные цепи подачи и развязки опущены. Энергию отдаёт источник накачки. Для вырожденного прибора достаточно одного контура с параметром, меняющимся около $2\omega_0$.
% Источник схемы: balakshy, PDF-с.9 / лист 7, рисунок вверху страницы.
% Белокопытов II, PDF-с.78–80, подтверждает трёхмодовое взаимодействие,
% но не является источником этой варикапной электрической схемы.

\paragraph{Вырожденное усиление.}
При $\omega_s=\omega_p/2$ на точной настройке повёрнутые квадратуры удовлетворяют $\dot X=-(\gamma-\kappa)X+F_X$, $\dot Y=-(\gamma+\kappa)Y+F_Y$. Ниже порога $\kappa<\gamma$ одна квадратура теряет меньше энергии, другая --- больше. Относительно режима без накачки стационарные \emph{амплитудные} коэффициенты
\begin{equation}
 G_+=\frac1{1-\kappa/\gamma},\qquad G_-=\frac1{1+\kappa/\gamma}.
\end{equation}
Усиление зависит от фазы сигнала. При $\kappa\to\gamma$ линейная модель расходится: начинается самовозбуждение с порогом \eqref{eq:paramthreshold}, а не бесконечная выходная мощность.

\paragraph{Невырожденный прибор.}
Две моды выбирают сигнал и холостую частоту $\omega_i=\omega_p-\omega_s$. Смешение «сигнал $\to$ холостая $\to$ сигнал» даёт сигналу добавочную энергию. При резонансе и неистощаемой накачке нормированные огибающие:
\begin{equation}\label{eq:twomodeparam}
 \dot a_s=-\gamma_sa_s+ga_i^*+F_s,\qquad
 \dot a_i^*=-\gamma_ia_i^*+g^*a_s.
\end{equation}
$g$ --- скорость параметрической связи, пропорциональная накачке и перекрытию мод; $\gamma_j$ --- декременты амплитуд. Исключив стационарную $a_i^*$, получаем $a_s=F_s/(\gamma_s-|g|^2/\gamma_i)$. Порог без входного сигнала:
\begin{equation}
 |g|^2>\gamma_s\gamma_i.
\end{equation}
Расстройка повышает порог; конкретный коэффициент через глубину модуляции зависит от схемы. При отсутствии входа в холостой порт его фаза подстраивается, поэтому усиление одного сигнала не зависит от его абсолютной фазы. Внешний портовый коэффициент дополнительно зависит от связи и согласования.

Для идеального реактивного трёхволнового взаимодействия соотношения Мэнли--Роу выражают равенство скоростей рождения пар:
\begin{equation}
 \frac{\dot W_s}{\omega_s}=\frac{\dot W_i}{\omega_i}
 =-\frac{\dot W_p}{\omega_p},\qquad \omega_p=\omega_s+\omega_i.
\end{equation}
Здесь только вклады взаимодействия; потери добавляют отдельно. В оптике аналогичен кристалл с $\chi^{(2)}$ и фазовым синхронизмом. Реактивный элемент позволяет малошумящее усиление, но не устраняет шумы потерь и квантовые ограничения. Конечную генерацию выше порога ограничивают насыщение, расстройка или истощение накачки.
